\documentclass[10pt]{beamer}
\usetheme{metropolis}
\usepackage{graphicx}
\usepackage{listings}
\usepackage{booktabs}
\usepackage{xcolor}
\usepackage{hyperref}

\definecolor{codebg}{HTML}{F5F5F5}
\definecolor{codegreen}{HTML}{2E7D32}
\definecolor{codegray}{HTML}{757575}
\definecolor{codeblue}{HTML}{1565C0}

\lstdefinestyle{pythonstyle}{
  language=Python, backgroundcolor=\color{codebg}, basicstyle=\ttfamily\scriptsize,
  keywordstyle=\color{codeblue}\bfseries, stringstyle=\color{codegreen},
  commentstyle=\color{codegray}\itshape, breaklines=true, frame=single,
  rulecolor=\color{codegray!30}, showstringspaces=false, tabsize=4,
  morekeywords={import,from,as,pd,plt,sns,np,fig,ax,ggplot,aes,
    geom_point,geom_line,geom_col,geom_histogram,geom_boxplot,geom_smooth,
    geom_density,geom_violin,geom_bar,geom_jitter,
    scale_color_brewer,scale_fill_manual,facet_wrap,facet_grid,
    theme_minimal,theme,labs,coord_flip,
    scatterplot,boxplot,violinplot,histplot,kdeplot,heatmap,
    regplot,lmplot,catplot,relplot,displot,pairplot,
    FacetGrid,set_theme,set_style,set_context}
}
\lstdefinestyle{rstyle}{
  language=R, backgroundcolor=\color{codebg}, basicstyle=\ttfamily\scriptsize,
  keywordstyle=\color{codeblue}\bfseries, stringstyle=\color{codegreen},
  commentstyle=\color{codegray}\itshape, breaklines=true, frame=single,
  rulecolor=\color{codegray!30}, showstringspaces=false, tabsize=2,
  morekeywords={library,ggplot,aes,geom_point,geom_smooth,geom_histogram,
    geom_boxplot,facet_wrap,theme_minimal,labs,scale_color_brewer}
}

\title{seaborn \& plotnine for Python}
\subtitle{Workshop 2 --- Module 7}
\author{Prof.Asoc.\ Endri Raco}
\institute{Data Visualization for Data Scientists \\ UNYT Tirana}
\date{2025/26}

\begin{document}

\begin{frame}
  \titlepage
\end{frame}

\begin{frame}{Module Roadmap}
  \begin{enumerate}
    \item Three Python viz libraries: matplotlib, seaborn, plotnine
    \item seaborn: figure-level vs axes-level functions
    \item seaborn gallery: scatter, box, violin, histogram, heatmap, pairplot
    \item plotnine: the ggplot2 clone for Python
    \item Decision tree: which library for which task
    \item R $\leftrightarrow$ Python Rosetta Stone
  \end{enumerate}
  \begin{exampleblock}{Learning Outcome}
    You will choose the right Python viz library for each task and
    translate any ggplot2 chart into seaborn or plotnine with minimal
    code changes.
  \end{exampleblock}
\end{frame}

\section{Three Python Libraries}

\begin{frame}{matplotlib vs seaborn vs plotnine}
  \begin{center}
    \includegraphics[width=0.95\textwidth]{figures_w02m07/three_libraries}
  \end{center}
\end{frame}

\begin{frame}{When to Use Each}
  \begin{center}
    \includegraphics[width=0.85\textwidth]{figures_w02m07/decision_tree}
  \end{center}
  \begin{alertblock}{Pro-Tip}
    If you are already fluent in ggplot2 (R), start with \textbf{plotnine}
    --- the syntax is identical. If you want built-in statistical aggregation
    and faceting, use \textbf{seaborn}. Fall back to \textbf{matplotlib}
    only for full manual control or highly customised layouts.
  \end{alertblock}
\end{frame}

\section{seaborn}

\begin{frame}{Function Taxonomy: Figure-Level vs Axes-Level}
  \begin{center}
    \includegraphics[width=0.92\textwidth]{figures_w02m07/sns_function_types}
  \end{center}
\end{frame}

\begin{frame}[fragile]{Figure-Level: Built-In Faceting}
\begin{lstlisting}[style=pythonstyle]
import seaborn as sns

# Figure-level: creates its own Figure + facets automatically
g = sns.relplot(
    data=apps, x="Reviews", y="Rating",
    hue="Type",                  # colour aesthetic
    col="Content_Rating",        # facet by column
    col_wrap=3,                  # wrap after 3 columns
    kind="scatter",              # scatter (or "line")
    height=3, aspect=1.2,        # panel size
    alpha=0.3, s=15)
g.set(xscale="log")             # log scale on all panels
g.set_titles("{col_name}")      # strip label format
g.tight_layout()
g.savefig("faceted_scatter.png", dpi=300)
\end{lstlisting}
  \begin{exampleblock}{Data Insight}
    Figure-level functions (\texttt{relplot}, \texttt{catplot}, \texttt{displot})
    create the entire figure and handle faceting with \texttt{col=} and
    \texttt{row=}. Axes-level functions (\texttt{scatterplot}, \texttt{boxplot})
    plot onto a single existing axes --- use them inside \texttt{plt.subplots()}.
  \end{exampleblock}
\end{frame}

\begin{frame}[fragile]{Axes-Level: Embed in Custom Layouts}
\begin{lstlisting}[style=pythonstyle]
# Axes-level: plot onto an existing axes
fig, axes = plt.subplots(1, 3, figsize=(12, 4), sharey=True)

# Panel 1: scatter
sns.scatterplot(data=df, x="x", y="y", hue="group", ax=axes[0])
axes[0].set_title("Scatter")

# Panel 2: box
sns.boxplot(data=df, x="group", y="y", palette="Set2", ax=axes[1])
axes[1].set_title("Box")

# Panel 3: violin
sns.violinplot(data=df, x="group", y="y", palette="Set2", ax=axes[2])
axes[2].set_title("Violin")

plt.tight_layout()
\end{lstlisting}
  \begin{alertblock}{Pro-Tip}
    The \texttt{ax=} parameter is the key: it tells seaborn to draw on
    your axes instead of creating a new figure. This is how you combine
    seaborn with matplotlib's \texttt{GridSpec} for custom layouts.
  \end{alertblock}
\end{frame}

\begin{frame}{seaborn Gallery: Six Chart Types}
  \begin{center}
    \includegraphics[width=0.95\textwidth]{figures_w02m07/sns_gallery}
  \end{center}
\end{frame}

\begin{frame}[fragile]{seaborn: Multi-Aesthetic Encoding}
  \begin{columns}[T]
    \begin{column}{0.52\textwidth}
\begin{lstlisting}[style=pythonstyle]
sns.scatterplot(
    data=df,
    x="gdp",
    y="life_exp",
    hue="continent",    # colour
    style="income",     # shape
    size="population",  # size
    sizes=(20, 300),    # size range
    alpha=0.6,
    palette="Set2")
\end{lstlisting}
      \vspace{4pt}
      {\small seaborn's \texttt{hue}, \texttt{style},
      and \texttt{size} map directly to ggplot2's
      \texttt{aes(color=, shape=, size=)}. Legends
      are generated automatically.}
    \end{column}
    \begin{column}{0.45\textwidth}
      \includegraphics[width=\textwidth]{figures_w02m07/sns_multi_aes}
    \end{column}
  \end{columns}
\end{frame}

\begin{frame}[fragile]{seaborn: Statistical Aggregation}
\begin{lstlisting}[style=pythonstyle]
# catplot with built-in mean + CI (like stat_summary in ggplot2)
g = sns.catplot(
    data=apps, x="Category", y="Rating",
    kind="bar",        # computes mean + CI automatically
    ci=95,             # 95% confidence interval
    palette="Set2",
    height=5, aspect=1.5)
g.set_xticklabels(rotation=45, ha="right")

# regplot: scatter + regression line + CI (like geom_smooth)
sns.regplot(data=df, x="x", y="y",
    ci=95, scatter_kws={"alpha": 0.3, "s": 15},
    line_kws={"color": "#E53935"})

# lmplot: regplot + faceting
sns.lmplot(data=df, x="x", y="y",
    hue="group", col="region", col_wrap=3,
    ci=95, height=3, aspect=1.2)
\end{lstlisting}
\end{frame}

\begin{frame}{seaborn FacetGrid: Built-In Faceting}
  \begin{center}
    \includegraphics[width=0.95\textwidth]{figures_w02m07/sns_facetgrid}
  \end{center}
\end{frame}

\section{plotnine}

\begin{frame}{ggplot2 $\rightarrow$ plotnine: Line-by-Line}
  \begin{center}
    \includegraphics[width=0.92\textwidth]{figures_w02m07/plotnine_syntax}
  \end{center}
\end{frame}

\begin{frame}[fragile]{plotnine: Complete Example}
  \begin{columns}[T]
    \begin{column}{0.52\textwidth}
\begin{lstlisting}[style=pythonstyle]
from plotnine import *
import pandas as pd

apps = pd.read_csv(
    "googleplaystore.csv")

(ggplot(apps.dropna(
        subset=["Rating"]),
    aes(x="Reviews", y="Rating",
        color="Type"))

 + geom_point(alpha=0.3, size=0.8)

 + geom_smooth(method="lm",
     se=True, color="grey")

 + scale_x_log10()

 + scale_color_manual(
     values={"Free": "#1565C0",
             "Paid": "#E53935"})

 + facet_wrap("~Content_Rating",
     ncol=2)

 + theme_minimal()

 + labs(title="Reviews vs Rating",
     x="Reviews (log)",
     y="Rating")
).save("chart.pdf",
    width=8, height=6, dpi=300)
\end{lstlisting}
    \end{column}
    \begin{column}{0.45\textwidth}
      \includegraphics[width=\textwidth]{figures_w02m07/plotnine_output}

      \vspace{4pt}
      {\small Every ggplot2 feature has
      a plotnine equivalent:
      \texttt{aes()}, \texttt{geom\_*()},
      \texttt{scale\_*()}, \texttt{facet\_*()},
      \texttt{theme\_*()}, \texttt{labs()}.
      The only differences are strings
      for column names and \texttt{+}
      at line start.}
    \end{column}
  \end{columns}
\end{frame}

\begin{frame}[fragile]{plotnine vs ggplot2: The Three Differences}
  \centering
  \begin{tabular}{lll}
    \toprule
    \textbf{Feature} & \textbf{R (ggplot2)} & \textbf{Python (plotnine)} \\
    \midrule
    Column names & Bare symbols: \texttt{aes(x = displ)} & Strings: \texttt{aes(x="displ")} \\
    Line continuation & \texttt{+} at end of line & \texttt{+} at start, inside \texttt{()} \\
    Save & \texttt{ggsave("p.pdf")} & \texttt{p.save("p.pdf")} \\
    \bottomrule
  \end{tabular}

  \vspace{8pt}
  \begin{exampleblock}{Data Insight}
    Everything else --- geom names, scale functions, facet syntax, theme
    elements --- is character-for-character identical. If you can write
    ggplot2, you can write plotnine in under 5 minutes.
  \end{exampleblock}
\end{frame}

\section{Rosetta Stone}

\begin{frame}{R $\leftrightarrow$ Python Translation Table}
  \begin{center}
    \includegraphics[width=0.95\textwidth]{figures_w02m07/rosetta}
  \end{center}
\end{frame}

\begin{frame}[fragile]{Same Chart, Three Ways}
  \begin{columns}[T]
    \begin{column}{0.32\textwidth}
      \textbf{R (ggplot2)}
\begin{lstlisting}[style=rstyle]
ggplot(df, aes(
    x = displ, y = hwy,
    color = class)) +
  geom_point(alpha=0.6) +
  geom_smooth(
    method = "lm") +
  scale_color_brewer(
    palette = "Set2") +
  facet_wrap(~cyl) +
  theme_minimal()
\end{lstlisting}
    \end{column}
    \begin{column}{0.32\textwidth}
      \textbf{Python (plotnine)}
\begin{lstlisting}[style=pythonstyle]
(ggplot(df, aes(
    x="displ", y="hwy",
    color="class"))
 + geom_point(alpha=0.6)
 + geom_smooth(
     method="lm")
 + scale_color_brewer(
     type="qual",
     palette="Set2")
 + facet_wrap("~cyl")
 + theme_minimal())
\end{lstlisting}
    \end{column}
    \begin{column}{0.32\textwidth}
      \textbf{Python (seaborn)}
\begin{lstlisting}[style=pythonstyle]
g = sns.lmplot(
    data=df,
    x="displ",
    y="hwy",
    hue="class",
    col="cyl",
    col_wrap=2,
    palette="Set2",
    height=3,
    aspect=1.2,
    ci=95,
    scatter_kws={
      "alpha": 0.6})
\end{lstlisting}
    \end{column}
  \end{columns}
\end{frame}

\section{Synthesis}

\begin{frame}{Library Selection Summary}
  \centering
  \begin{tabular}{lp{6cm}}
    \toprule
    \textbf{Use Case} & \textbf{Recommended Library} \\
    \midrule
    Want ggplot2 syntax in Python & \texttt{plotnine} \\
    Statistical charts with auto-aggregation & \texttt{seaborn} \\
    Faceted statistical charts & \texttt{seaborn} (figure-level) or \texttt{plotnine} \\
    Custom multi-panel dashboard & \texttt{matplotlib} + \texttt{GridSpec} \\
    Embed seaborn in custom layout & \texttt{seaborn} axes-level + \texttt{ax=} \\
    Heatmaps, pairplots, clusterplots & \texttt{seaborn} \\
    Maximum customisation & \texttt{matplotlib} \\
    \bottomrule
  \end{tabular}
\end{frame}

\begin{frame}{Key Takeaways}
  \begin{enumerate}
    \item \textbf{matplotlib}: low-level, full control, verbose.
          \textbf{seaborn}: statistical, auto-aggregation, faceting.
          \textbf{plotnine}: ggplot2 clone, grammar syntax.
    \item seaborn has two API levels: \textbf{figure-level} (creates figure,
          has faceting) and \textbf{axes-level} (draws on existing ax).
    \item seaborn's \textbf{hue/style/size} map to ggplot2's
          \textbf{color/shape/size} aesthetics.
    \item plotnine syntax is \textbf{identical} to ggplot2 except: strings
          for column names, \texttt{+} at line start, \texttt{.save()}.
    \item The \textbf{Rosetta Stone} table translates any ggplot2 chart
          to seaborn or matplotlib with known equivalences.
  \end{enumerate}
\end{frame}

\begin{frame}{Essential Readings}
  \begin{itemize}
    \item seaborn tutorial: \texttt{https://seaborn.pydata.org/tutorial.html}
    \item plotnine documentation: \texttt{https://plotnine.org}
    \item VanderPlas, J. (2016). \emph{Python Data Science Handbook},
          Chapter 4 (Visualization with Matplotlib).
    \item Waskom, M. (2021). ``seaborn: statistical data visualization.''
          \emph{JOSS}, 6(60), 3021.
  \end{itemize}
  \vspace{6pt}
  \textbf{Next Module}: labs(), Annotations \& Storytelling Layers (W02-M08)
\end{frame}

\end{document}
